\chapter{Intro}
\label{chap:2}
\begin{center}
    Варіант $ = \max \left\{5 ; 10\right\} = 10$
\end{center}

\noindent Мета комп'ютерного практикуму: Практична реалізація статистичної атаки на комбінувальний генератор гами; 
набуття навичок роботи з системами комп'ютерної алгебри в SAGE.

\noindent Постановка задачі:
\begin{enumerate}
    \item Обчислити коефіцієнти Уолша-Адамара функції $f$. За допомогою цих коефіцієнтів знайти \underline{всі афінні}
        статистичні аналоги $f$ та ймовірності збігу цих аналогів з $f$.
    \item Вибрати набір статистичних аналогів $g_{1}, \dots, g_{r}$, що будуть використані для відновлення початкового
        стану генератора. Для всіх функцій з цього набору обчислити необхідну кількість матеріалу $T_{1}, \dots, T_{r}$
        відповідно, зафіксувавши $\delta = 0.01$. Побудувати низку атак на початкові стани регістрів, тим самим 
        повністю відновивши початковий стан генератора гами.
    \item Перевірити, що початковий стан генератора відновлено правильно, шляхом генерування відрізка гами відповідної
        довжини й порівняти його з вхідними даними.
\end{enumerate}

\section{Вхідні дані}
\begin{enumerate}
    \item Шість лінійних регістрів зсуву довжини $m = 10$ з поліномами зворотного зв’язку є наступними:
        \begin{align*}
                & p_{1}(x) = x_{10} \oplus x_{3} \oplus 1, \quad p_{2}(x) = x_{10} \oplus x_{7} \oplus 1,                                                     \\
                & p_{3}(x) = x_{10} \oplus x_{4} \oplus x_{3} \oplus x_{1} \oplus 1, \quad p_{4}(x) = x_{10} \oplus x_{8} \oplus x_{3} \oplus x_{2} \oplus 1, \\
                & p_{5}(x) = x_{10} \oplus x_{8} \oplus x_{4} \oplus x_{3} \oplus 1, \quad p_{6}(x) = x_{10} \oplus x_{8} \oplus x_{5} \oplus x_{1} \oplus 1. \\
        \end{align*}
    \item Комбінувальна функція $f(x_{0}, \ldots, x_{63})$ генератора гами у файлі \textbf{combinfuncV10.txt}
        \begin{equation*}
            x_{0} x_{2} x_{4} + x_{1} x_{2} x_{4} + x_{1} x_{4} + x_{2} x_{3} x_{4} + x_{2} x_{4} x_{5} + x_{2} x_{4} + x_{2} + x_{3} x_{4} + x_{4} + 1
        \end{equation*}
        $\deg f = 3$, вага булевої функції -- кількість входів, на яких вона дорівнює 1 за таблицею істинності. 
        В нашому випадку $\mathrm{wt}(f) = 32$ (збалансована).
    \item Відрізок гами $\gamma_{1}, \ldots, \gamma_{N}$ довжини $N = 50000$ бітів у файлі \textbf{gamma2V10.txt}.

\end{enumerate}

\section{Хід виконання роботи}

Використовував пакет комп'ютерної алгебри SageMath (версія~10.7) на WSL~2 Ubuntu, встановлений при виконанні
попередньої лаби.

\subsection{Коефіцієнти Уолша-Адамара та афінні статистичні аналоги}

Перетворення Уолша-Адамара~\cite{siegenthaler1985} є основним інструментом аналізу кореляційних властивостей булевих 
функцій. Для функції $f \in B_n$ коефіцієнт Уолша-Адамара в точці $\alpha \in V_n$ визначається як:
\begin{equation*}
    \hat{f}(\alpha) = \sum_{x \in V_n} (-1)^{f(x) \oplus \langle \alpha, x \rangle},
\end{equation*}
де $\langle \alpha, x \rangle = \bigoplus\limits_{i=0}^{n-1} \alpha_i x_i$~ -- скалярний добуток у $\mathrm{GF}(2)$.

\noindent Кожен ненульовий коефіцієнт $\hat{f}(\alpha)$ визначає афінний статистичний аналог $g = \langle \alpha, x \rangle \oplus c$, 
де $c = 0$ при $\hat{f}(\alpha) > 0$ та $c = 1$ при $\hat{f}(\alpha) < 0$. В даному випадку $c$ виступає аналогом 
константи (вільним членом) в афінному представленні. Зміщення (bias) аналога:
\begin{equation*}
    \varepsilon = \frac{|\hat{f}(\alpha)|}{2^n},
\end{equation*}
а ймовірність збігу $f$ та $g$:
\begin{equation*}
    \Pr \left(f(X) = \underbrace{\alpha_{1} X_{n} \oplus \ldots \alpha_{n} X_{n} \oplus c}_{g(X)}\right) = 
    \frac{1}{2}\left(1 + (-1)^{c} \cdot \varepsilon\right).
\end{equation*}

\begin{figure}[ht]
    \centering
    \includegraphics[width=0.75\linewidth]{screenshots/step1.png}
    \caption{Крок 1: запуск і комбінувальна функція}
\end{figure}

\noindent Коефіцієнти обраховували вбудованою функцією SAGE \texttt{walsh\_hadamard\_transform()}.

\newpage
\begin{figure}[ht]
    \centering
    \begin{minipage}{.48\textwidth}
        \centering
        \includegraphics[width=.75\linewidth]{screenshots/step2.1.png}
        \captionof{figure}{Крок 2: коефіцієнти \\ Уолша-Адамара}
    \end{minipage}
    \begin{minipage}{.48\textwidth}
        \centering
        \includegraphics[width=.8\linewidth]{screenshots/step2.2.png}
        \captionof{figure}{Крок 2: коефіцієнти \\ Уолша-Адамара (continue)}
    \end{minipage}
\end{figure}

\noindent Знайдено $10$ ненульових коефіцієнтів:

\begin{figure}[ht]
    \centering
    \includegraphics[width=0.75\linewidth, scale=0.85]{screenshots/step3.png}
    \caption{Крок 3: $g$ аналоги і $\mathrm{cor}(f)$}
\end{figure}

\noindent Кореляційна імунність~\cite{siegenthaler1985}: $\mathrm{cor}(f) = 0$, оскільки існує ненульовий коефіцієнт 
ваги~$1$ ($\hat{f}(001000) = -32$, що відповідає $g = x_2 \oplus 1$). Це означає наявність аналога від \emph{одного} регістру. 
На даному screenshot (і далі) під "+"{} мається на увазі $\oplus$.

\subsection{Вибір аналогів та необхідна кількість матеріалу}

Атака Зігенталера~\cite{siegenthaler1985} використовує метод максимуму правдоподібності: 
для кожного набору кандидатів на початкові стани $k$ регістрів генерується відрізок гами довжини $T$ через 
аналог $g$ і обчислюється відстань Гемінга до відповідного відрізку вхідної гами. Набір з мінімальною такою 
відстанню є найбільш імовірним початковим станом гами.

\noindent Необхідна кількість матеріалу~\cite{carlet2010} (та й в наших лекціях це було):
\begin{equation*}
    T = \left\lceil 8 \cdot \varepsilon^{-2} \cdot \ln\bigl(2^{mk} / \delta\bigr) \right\rceil,
\end{equation*}
де $m = 10$ (довжина регістру), $k$ --- кількість регістрів в аналозі, $\delta = 0.01$ -- фіксоване.

\begin{figure}[ht]
    \centering
    \includegraphics[width=0.75\linewidth, scale=0.6]{screenshots/step5.png}
    \caption{Кроки 5--7: зчитування гами, параметри ЛРЗЗ, обчислення $T$}
\end{figure}

\noindent Стратегія вибору аналогів визначається автоматично за "жадібним"{} принципом:
\begin{itemize}[noitemsep]
    \item Аналоги $g$ сортуються за зростанням $k$ (менше регістрів~--- менше кандидатів для перебору) та за 
        спаданням~$\varepsilon$ (більше зміщення --- менше потрібно матеріалу $T$).
    \item На кожному етапі обирається найдешевший аналог, що покриває лише ще невідновлені регістри.
    \item Якщо для якогось регістру не залишилось "чистого"{} аналога (всі аналоги, що його містять, 
        мають хоча б один уже відновлений регістр), цей регістр відновлюватиметься повним перебором.
\end{itemize}

\noindent Для нашого варіанту ця процедура дала наступну послідовність атак:

\begin{enumerate}
    \item $g_1 = x_2 \oplus 1$ ($k=1$, $\varepsilon = 0.5$, регістр R3):
    $T_1 = 370$ бітів, $1023$ кандидати.
    \item $g_2 = x_1 \oplus x_3$ ($k=2$, $\varepsilon = 0.25$, регістри R2, R4):
    $T_2 = 2364$ бітів, $1\,046\,529$ кандидатів.
    \item $g_3 = x_0 \oplus x_5$ ($k=2$, $\varepsilon = 0.25$, регістри R1, R6):
    $T_3 = 2364$ бітів, $1\,046\,529$ кандидатів.
    \item Регістр R5 відновлюється повним перебором ($1023$ кандидати).
\end{enumerate}

\begin{remark}
    R5 потрапив у повний перебір, оскільки єдиний аналог з $k \leq 2$, що його містить~--- $g = x_2 \oplus x_4 \oplus 1$ 
    (R3, R5) -- має R3 вже відновленим на першому етапі, а написаний скрипт пропускає аналоги з частково відомими 
    регістрами, бо я подумав що, при $m = 10$ перебрати $1023$ кандидатів швидше, ніж ускладнювати логіку для 
    часткової атаки, хех. І це справді виявилося так:) (див. час роботи нижче)
\end{remark}

\begin{remark}
    Ще варто зазначити, що якби функція $f$ мала вищу кореляційну імунність (наприклад, $\mathrm{cor}(f) = 2$), то 
    аналогів від одного регістру не існувало б, і той самий алгоритм автоматично почав би з аналогів при $k \geq 3$.
\end{remark}

\subsection{Оцінка складності}

Згідно з~\cite{siegenthaler1985} (і конспектом), часова складність атаки з аналогом від~$k$ регістрів становить:
\begin{equation*}
    O \left(T \cdot 2^{mk} + nm \cdot 2^{m(n-k)}\right),
\end{equation*}
де перший доданок -- складність статистичної фази (перебір $2^{mk}$ кандидатів, для кожного порівняння $T$ бітів), 
а другий -- складність повного перебору залишкових $n - k$ регістрів. Для нашого варіанту ($n = 6$, $m = 10$):

\begin{center}
    \begin{tabular}{|l|c|c|c|c|}
        \hline
        \textbf{Етап}          & $k$                                         & $T \cdot 2^{mk}$                           & $nm \cdot 2^{m(n-k)}$ & \textbf{Сумарно} \\
        \hline
        $g_1 = x_2 \oplus 1$   & 1                                           & $370 \cdot 2^{10} \approx 3.8 \cdot 10^5$
                               & $60 \cdot 2^{50} \approx 6.8 \cdot 10^{16}$
                               & $\approx 6.8 \cdot 10^{16}$                                                                                                         \\
        $g_2 = x_1 \oplus x_3$ & 2                                           & $2364 \cdot 2^{20} \approx 2.5 \cdot 10^9$
                               & $60 \cdot 2^{40} \approx 6.6 \cdot 10^{13}$
                               & $\approx 6.6 \cdot 10^{13}$                                                                                                         \\
        $g_3 = x_0 \oplus x_5$ & 2                                           & $2364 \cdot 2^{20} \approx 2.5 \cdot 10^9$
                               & $60 \cdot 2^{40} \approx 6.6 \cdot 10^{13}$
                               & $\approx 6.6 \cdot 10^{13}$                                                                                                         \\
        \hline
    \end{tabular}
\end{center}

Однак атаки виконуються \underline{послідовно}: після етапу~1 відомий стан R3, тому на етапі~2 перебираються лише 
$2$ невідомих регістри (а не $n - k = 4$ залишкових повним перебором). Аналогічно після етапів~1--3 залишається 
лише R5, для якого перебір $2^{10} - 1 = 1023$ кандидатів є тривіальним.

\noindent Фактична складність трьох статистичних фаз:
\begin{equation}
    \label{eq:summary_complexity}
    T_1 \cdot 2^{m} + T_2 \cdot 2^{2m} + T_3 \cdot 2^{2m} + 2^{m}
    \approx 3.8 \cdot 10^5 + 2.5 \cdot 10^9 + 2.5 \cdot 10^9 + 10^3 
    \approx 5 \cdot 10^9,
\end{equation}
що на багато порядків менше за повний перебір $2^{nm} = 2^{60} \approx 1.15 \cdot 10^{18}$.

\subsection{Передобчислення послідовностей ЛРЗЗ}

Перед безпосереднім запуском атак виконується передобчислення: для кожного з $6$ регістрів і кожного з $2^m - 1 = 1023$ 
можливих ненульових початкових станів  -- генерується вихідна послідовність ЛРЗЗ довжиною~$T_{\max}$ 
(максимальне~$T$ серед усіх обраних аналогів).

Кожна послідовність зберігається як одне велике ціле число (\texttt{big integer}), де біт~$i$ дорівнює виходу ЛРЗЗ 
у момент часу~$i$. Всього обчислюється $6 \times 1023 = 6138$ послідовностей.

Це дозволяє обчислювати відстань Гемінга між гамою та виходом аналога за допомогою  лише двох побітових операцій: 
XOR двох великих чисел та підрахунок одиничних бітів (\texttt{popcount}). Без передобчислення довелося б на кожному 
з мільйонів кандидатів заново генерувати послідовність ЛРЗЗ побітово, що очевидно було б повільніше (тому я навіть 
цим і не заморочувався). Час передобчислення склав $\approx 20$~с.

\begin{figure}[ht]
    \centering
    \includegraphics[width=0.85\linewidth, scale=0.8]{screenshots/step8.png}
    \caption{Крок 8: передобчислення послідовностей LFSR}
\end{figure}

\subsection{Результати статистичних атак}

\begin{enumerate}
    \item[] \textbf{Етап 1: $g_1 = x_2 \oplus 1$ (регістр R3).}
        Перебір $1023$ кандидатів, $T_1 = 370$ бітів.
        Мінімальне значення статистики (абсолютне): $84$, відносне: $84/370 = 0.2270$.
        Друге найкраще (абсолютне): $156$, відносне: $156/370 = 0.4216$.
        Час: $< 0.01$~с. \textbf{R3 = 972} (\texttt{0b1111001100}).
    \item[] \textbf{Етап 2: $g_2 = x_1 \oplus x_3$ (регістри R2, R4).}
        Перебір $1\,046\,529$ кандидатів, $T_2 = 2364$ бітів.
        Мінімальне (абс.): $899$, відносне: $899/2364 = 0.3803$.
        Друге найкраще (абс.): $1012$, відносне: $1012/2364 = 0.4281$.
        Час: $\approx 11.4$~с. \textbf{R2 = 610} (\texttt{0b1001100010}),
        \textbf{R4 = 272} (\texttt{0b0100010000}).
    \item[] \textbf{Етап 3: $g_3 = x_0 \oplus x_5$ (регістри R1, R6).}
        Перебір $1\,046\,529$ кандидатів, $T_3 = 2364$ бітів.
        Мінімальне (абс.): $873$, відносне: $873/2364 = 0.3693$.
        Друге найкраще (абс.): $998$, відносне: $998/2364 = 0.4222$.
        Час: $\approx 11.2$~с. \textbf{R1 = 223} (\texttt{0b0011011111}),
        \textbf{R6 = 672} (\texttt{0b1010100000}).
    \item[] \textbf{Етап 4: повний перебір R5.}
        Перебір $1023$ кандидатів (із зафіксованими R1--R4, R6), $T = 2000$ бітів.
        Мінімальне (абс.): $0$ (точний збіг).
        Час: $\approx 12.9$~с. \textbf{R5 = 146} (\texttt{0b0010010010}).
\end{enumerate}

\begin{figure}[ht]
    \centering
    \includegraphics[width=0.85\linewidth]{screenshots/step9.png}
    \caption{Кроки 9--10: статистичні атаки, повний перебір R5}
\end{figure}

\begin{otherremark}
    При повному переборі R5 ми вже знаємо стани всіх інших 5 регістрів. Тому для кожного кандидата на R5 генеруємо 
    гаму через повну комбінувальну функцію $f$ (а не через наближений аналог $g$). Коли кандидат правильний, тобто 
    згенерована гама точно збігається з вхідною, тому відстань Гемінга = 0.
\end{otherremark}

\noindent Отже, відновлений початковий стан генератора має вид:

\begin{center}
    \begin{tabular}{|c|c|c|}
        \hline
        \textbf{Регістр} & \textbf{decimal} & \textbf{binary}     \\
        \hline
        R1               & 223              & \texttt{0011011111} \\
        R2               & 610              & \texttt{1001100010} \\
        R3               & 972              & \texttt{1111001100} \\
        R4               & 272              & \texttt{0100010000} \\
        R5               & 146              & \texttt{0010010010} \\
        R6               & 672              & \texttt{1010100000} \\
        \hline
    \end{tabular}
\end{center}

\subsection{Перевірка}

Зі знайдених початкових станів було згенеровано $N = 50000$ бітів гами за допомогою повної комбінувальної 
функції $f$ та порівняно з вхідними даними.

Результат як очікувано -- усі $50000$ з $50000$ бітів збіглися (\textbf{100\%}) -- перемога!.

\begin{figure}[ht]
    \centering
    \includegraphics[width=0.85\linewidth]{screenshots/step11.png}
    \caption{Крок 11: верифікація та підсумок}
\end{figure}

\section{Висновки}

У ході роботи ми реалізовували статистичну атаку Зігенталера на комбінувальний генератор гами, де $n = 6$ 
регістрів довжини $m = 10$, $\deg f = 3$
\begin{enumerate}
    \item Обчислили коефіцієнти Уолша-Адамара функції~$f$. Знайдено $10$ афінних статистичних аналогів. 
        Кореляційна імунність склала $\mathrm{cor}(f) = 0$, що означає наявність аналога від одного регістру.
    \item Для відновлення станів використано 3 аналоги ($k = 1, 2, 2$) та повний перебір для одного залишкового 
        регістру. Загальний час атаки: $< 0.01 + 11.4 + 11.2 + 12.9 \approx 36.5$~с (без передобчислення).
    \item Відновлено початкові стани всіх 6 регістрів:
        \begin{equation*}
            R1 = 223, \, R2 = 610, \, R3 = 972, \, R4 = 272, \, R5 = 146, \, R6 =672.
        \end{equation*}
    \item Початковий стан перевірено шляхом генерування гами, всі $N = 50000$ бітів збіглися з вхідними даними.
\end{enumerate}

Низька кореляційна імунність ($\mathrm{cor}(f) = 0$) зробила атаку надзвичайно ефективною: наявність аналога від одного 
регістру дозволила відновити R3 за частки секунди. Загальна складність атаки~$\approx 5 \cdot 10^9$ операцій~\eqref{eq:summary_complexity}, 
що на $9$ порядків менше за повний перебір $2^{60} \approx 1.15 \cdot 10^{18}$ можливих станів генератора~\cite{carlet2010}.

\section{Code listing section}

\lstinputlisting[caption={Статистична атака на комбінувальний генератор Gamma з використанням SAGE}, label={lst:attack}]{attack2.sage}

